\chapter{GDL/GWL Version of this Tutorial}

This appendix describes the GDL version of this document. GDL is the open-source core
of Genworks\textregistered{} GDL, and is a significant superset of ANSI CL. The GDL version of this
tutorial is not currently available, but will be posted to

\begin{center}
\texttt{http://www.genworks.com/downloads}
\end{center}

when it becomes available. The GDL version of this document will be similar to the CL version, but will include examples using GDL and GWL\footnote{GWL stands for the Generative Web Language, an HTML templating system for GDL} (Generative Web Language)
in addition to base CL.

GDL (Generative Design Language) is a domain-specific language extension to Common
Lisp for knowledge-based engineering and product design applications. GDL objects extend
CLOS\index{CLOS} objects with features such as:

\begin{itemize}
\item Cached, demand-driven computation
\item Automatic dependency tracking
\item Tree-structured object hierarchies
\item Quantification (automatic generation of sequences of child objects)
\item Optional default values for slots
\end{itemize}

GWL extends GDL to provide a declarative system for generating web-based user interfaces.
GWL allows developers to specify web page layouts and interactions using GDL object
definitions, automatically handling the HTTP request/response cycle and session management.

For more information about GDL and GWL, visit:

\begin{center}
\texttt{http://www.genworks.com}
\end{center}

The open-source Gendl system can be found at:

\begin{center}
\texttt{http://www.gendl.org}
\end{center}
